\section*{Chapitre \ref{ch:b2:ore-deposits} --- Gisements de minerai}

\begin{solution}{exo:b2:ore-deposits:1}
(a) $\qty{0.5}{\percent} = \qty{5000}{ppm}$, et $5000/28 = 179$. (b) $\qty{6}{g/t} = \qty{6000}{ng/g}$,
et $6000/1.5 = 4000$. (c) La chromite contient $2 \times 52.00/(55.85 + 104.00 + 64.00) = 0.4646$
de chrome en masse~; la chromitite,
$0.9 \times 0.4646 = 0.418$, soit \qty{418000}{ppm}, et $418000/92 = 4540$.
\end{solution}

\begin{solution}{exo:b2:ore-deposits:2}
(a) $914 \times 0.0053 = \qty{4.84}{Mt}$ de cuivre et $914 \times 10^6 \times 0.40 = \num{3.66e8}$ g,
\qty{366}{t} d'or~; $1640 \times 0.0044 = \qty{7.22}{Mt}$ et \qty{738}{t}. (b) $1640 - 914 =
\qty{726}{Mt}$ contenant $7.216 - 4.844 = \qty{2.37}{Mt}$ de cuivre, soit une \omterm{def:b2:ore-deposits:grade}{teneur} moyenne de
$2.372/726 = \qty{0.327}{\percent}$~: les blocs supplémentaires se situent sous la coupure de \qty{0.3}{\percent}
et juste au-dessus d'elle. (c) Les deux estimations répondent à la question pour deux coupures~; aucune ne donne la taille
du \omterm{def:b2:ore-deposits:deposit}{gisement} sans sa coupure.
\end{solution}

\begin{solution}{exo:b2:ore-deposits:3}
(a) $(19.3 - 1.00)/(2.65 - 1.00) = 18.3/1.65 = 11.1$. (b) Des vitesses égales exigent $(\rho_g - 1)d_g^2 =
(\rho_q - 1)d_q^2$, donc $d_g/d_q = \sqrt{1.65/18.3} = 0.300$~: un grain d'or se dépose comme un grain de quartz
environ trois fois plus gros. (c) $d_q = 0.1\sqrt{3.7/1.65} = \qty{0.150}{mm}$.
\end{solution}

\begin{solution}{exo:b2:ore-deposits:4}
Par unité de surface et à densité égale, la couverture conserve son propre cuivre et reçoit celui de la colonne
lessivée~: $g = g_0(h_e + \eta h_l)/h_e$. (a) $0.4 \times (50 + 150)/50 = \qty{1.6}{\percent}$.
(b) $0.4 \times (50 + 0.8 \times 150)/50 = \qty{1.36}{\percent}$. (c) $h_l = 50 \times (1.7/0.4 - 1)
= \qty{163}{m}$. Ces \omterm{def:b2:ore-deposits:grade}{teneurs} se situent dans la fourchette de 0.4--\qty{1.7}{\percent} des couvertures
d'enrichissement des \omterm{def:b2:ore-deposits:deposit}{gisements} de porphyre.
\end{solution}

\begin{solution}{exo:b2:ore-deposits:5}
(a) Platine, $C_S = 10 \times 10^4 (R + 1)/(R + 10^4)$ ppb~: \qty{1.00}{ppm} pour $R = 100$, \qty{9.10}{ppm}
pour \num{1000}, \qty{50.0}{ppm} pour \num{10000}, \qty{90.9}{ppm} pour \num{100000}.
Nickel, $C_S = 200 \times 500 (R + 1)/(R + 500)$ ppm~: \qty{1.68}{\percent}, \qty{6.67}{\percent},
\qty{9.52}{\percent}, \qty{9.95}{\percent}. (b) $D\,C_L$~: \qty{100}{ppm} de platine,
\qty{10}{\percent} de nickel. (c) Le nickel est à moins de \qty{5}{\percent} de son plafond
à partir de $R = \num{10000}$, si bien que sa \omterm{def:b2:ore-deposits:grade}{teneur} varie à peine sur un facteur dix en $R$~;
le platine gagne encore presque le double à cet endroit, si bien que sa \omterm{def:b2:ore-deposits:grade}{teneur} mesure $R$.
\end{solution}

\begin{solution}{exo:b2:ore-deposits:6}
(a) $\mu_M = 2.373 - 0.364 - 2 = 0.009$~; $\sigma_M = 0.702$~; médiane $10^{0.009} = \qty{1.02}{Mt}$.
(b) $z = (1 - 0.009)/0.702 = 1.412$, $1 - \Phi(1.412) = 0.0790$~: environ
33 \omterm{def:b2:ore-deposits:deposit}{gisements} sur 422. (c) $\sigma = \sqrt{0.654^2 + 0.199^2} = 0.684$, $z = 1.449$,
$1 - \Phi = 0.0737$, soit environ 31~: la faible corrélation ajoute quelques géants.
\end{solution}

\begin{solution}{exo:b2:ore-deposits:7}
Moyenne de tous les blocs~: $0.40\,e^{0.18} = \qty{0.479}{\percent}$. Pour $c = 0.3$~: $z = \ln(0.40/0.3)/0.60 =
0.4795$, $\Phi(z) = 0.6842$, $\Phi(z + s) = 0.8598$~: $T = \qty{684}{Mt}$, $\bar g = 0.479 \times
0.8598/0.6842 = \qty{0.602}{\percent}$, $M = 1000 \times 0.00479 \times 0.8598 = \qty{4.12}{Mt}$.
Pour $c = 0.6$~: $z = -0.6758$, $\Phi = 0.2496$, $\Phi(z + s) = 0.4698$~: \qty{250}{Mt} à
\qty{0.901}{\percent}, \qty{2.25}{Mt}. Sur la figure, à une coupure de \qty{0.3}{\percent}, la courbe bleue
donne environ \qty{680}{Mt} et la courbe rouge environ \qty{0.6}{\percent}.
\end{solution}

\begin{solution}{exo:b2:ore-deposits:8}
(a) Fluide~: $\num{5e6}/(\num{1e-3} \times 0.5) = \qty{1.00e10}{t}$~; \omterm{def:b2:magma-emplacement:magma}{magma}~: $\num{1e10}/0.035 =
\qty{2.86e11}{t}$. (b) $\num{2.86e14}/2600 = \qty{1.10e11}{m^3}$, soit \qty{110}{km^3}, un cube de \omterm{def:b2:magma-emplacement:magma}{magma}
d'environ \qty{4.8}{km} de côté. (c) Le \omterm{def:b2:magma-emplacement:magma}{magma} contenait $\num{2.86e11}
\times \num{50e-6} = \qty{1.43e7}{t}$ de cuivre, et le minerai en renferme $5/14.3 = 0.350$ de celui-ci.
\end{solution}

\begin{solution}{exo:b2:ore-deposits:9}
(a) Logarithmes~: moyenne 1.411, écart-type 0.0664~; médiane $10^{1.411} = \qty{25.8}{ppm}$~;
seuil $10^{1.411 + 0.133} = 10^{1.544} = \qty{35.0}{ppm}$. (b) $(4 \times 400 + 36 \times 25.8)/40 =
\qty{63.2}{ppm}$, au-dessus de \qty{35.0}{ppm}~: toujours anomal après dilution
par neuf parts de sédiment propre. (c) $4 \times 400 + (A - 4) \times 25.8 = 35.0\,A$ donne $A = 4 \times (400 -
25.8)/(35.0 - 25.8) = \qty{163}{km^2}$.
\end{solution}

\begin{solution}{exo:b2:ore-deposits:10}
(a) $e^{su}\,e^{-u^2/2} = e^{-(u - s)^2/2 + s^2/2}$~; diviser par $\sqrt{2\pi}$. (b) $E[e^{\mu + su}] =
e^{\mu}\int e^{su}\varphi(u)\,\mathrm{d}u = e^{\mu + s^2/2}\int\varphi(u - s)\,\mathrm{d}u = e^{\mu +
s^2/2}$. (c) $\Phi$ croît et $s > 0$, donc $\Phi(z + s) > \Phi(z)$~: les blocs conservés sont plus riches que les blocs
écartés, si bien que chaque coupure conserve une part de métal plus grande que la part de roche~; l'écart
croît avec $s$, la variabilité du \omterm{def:b2:ore-deposits:deposit}{gisement}.
\end{solution}

\begin{solution}{exo:b2:ore-deposits:11}
(a) $D(R + 1)/(R + D) = fD$ donne $R + 1 = fR + fD$, $R = (fD - 1)/(1 - f)$. (b) $f = 0.5$~: $R = D - 2$, soit
298 et \num{9998}~; $f = 0.9$~: $R = 9D - 10$, soit \num{2690} et \num{89990}. (c) Le nickel atteint \qty{50}{\percent}
de son plafond dans l'intervalle des minerais nickel--cuivre (de 100 à \num{2000}), et 90 juste au-delà~; le
platine exige $R$ de \num{10000} à \num{100000}, l'intervalle des récifs à éléments du groupe du platine.
\end{solution}

\begin{solution}{exo:b2:ore-deposits:12}
(a) Avec $\ln M = m + s u$, les \omterm{def:b2:ore-deposits:deposit}{gisements} au-dessus du quantile ont $u > \Phi^{-1}(q)$~; leur métal vaut
$\int_{\Phi^{-1}(q)}^{\infty} e^{m + su}\varphi(u)\,\mathrm{d}u = e^{m + s^2/2}[1 - \Phi(\Phi^{-1}(q) -
s)]$, et en divisant par le total $e^{m + s^2/2}$ on obtient $\Phi(s - \Phi^{-1}(q))$. (b) $s = 0.702 \ln 10 =
1.616$~; $\Phi(1.616 - 1.282) = \Phi(0.334) = 0.631$. (c) $\Phi(1.616 - 2.326) = \Phi(-0.710) = 0.239$,
contre 0.266 pour les quatre plus grands~: les plus grands \omterm{def:b2:ore-deposits:deposit}{gisements} sont un peu plus gros que ne le prévoit la
loi log-normale.
\end{solution}

\begin{solution}{pb:b2:ore-deposits:1}
\textbf{1.} \qty{330}{m}~; $\sum l\,g = 138.4$, moyenne $138.4/330 = \qty{0.419}{\percent}$.
\textbf{2.} $3.24/8 = \qty{0.405}{\percent}$~; chaque intervalle représente une longueur de carotte, donc un
volume de roche : un intervalle long doit compter plus qu'un intervalle court.
\textbf{3.} $230/330 = 0.697$, à $118.6/230 = \qty{0.516}{\percent}$.
\textbf{4.} $4190/28 = 150$.
\textbf{5.} $8100/28 = 289$.
\textbf{6.} $0.38\,e^{0.55^2/2} = 0.38\,e^{0.1513} = \qty{0.442}{\percent}$~; le forage, à
\qty{0.419}{\percent}, est à moins de \qty{6}{\percent} de cette valeur : représentatif.
\textbf{7.} $\pi \times 450^2 \times 700 = \qty{4.45e8}{m^3}$~; $\times 2.7 = \num{1.20e9}$ t,
\qty{1200}{Mt}.
\textbf{8.} $z = \ln(0.38/0.3)/0.55 = 0.4298$~; $\Phi(z) = 0.6663$.
\textbf{9.} $z = 1.167$, 0.4298, $-0.4990$~; $\Phi = 0.8784$, 0.6663, 0.3089~; $T = \qty{1050}{Mt}$,
\qty{800}{Mt}, \qty{371}{Mt}.
\textbf{10.} $\Phi(z + s) = \Phi(0.9798) = 0.8364$~; $\bar g = 0.442 \times 0.8364/0.6663 =
\qty{0.555}{\percent}$.
\textbf{11.} $0.442 \times 0.9570/0.8784 = \qty{0.482}{\percent}$~; $0.442 \times 0.5203/0.3089 =
\qty{0.744}{\percent}$.
\textbf{12.} \omterm{def:b2:ore-deposits:grade}{Tonnage} $\times\, 370.7/799.6 = 0.464$~; grade $\times\, 0.744/0.555 = 1.34$.
\textbf{13.} À la médiane, \qty{0.38}{\percent} ($z = 0$)~: \qty{600}{Mt} à $0.442 \times
\Phi(0.55)/0.5 = 0.442 \times 0.7088/0.5 = \qty{0.627}{\percent}$.
\textbf{14.} $1200 \times 0.00442 = \qty{5.30}{Mt}$.
\textbf{15.} $(\log_{10} 1200 - 2.373)/0.654 = (3.0792 - 2.373)/0.654 = 1.080$~; $\Phi = 0.8599$~: le
86e centile, un \omterm{def:b2:ore-deposits:deposit}{gisement} de grande taille.
\textbf{16.} $(\log_{10} 0.442 + 0.364)/0.199 = (-0.3546 + 0.364)/0.199 = 0.047$~; $\Phi = 0.5187$~: une
\omterm{def:b2:ore-deposits:grade}{teneur} ordinaire, au 52e centile.
\textbf{17.} 0.702 avec la corrélation, 0.684 sans elle.
\textbf{18.} $\mu_M = 0.009$~; $(\log_{10} 5.30 - 0.009)/0.702 = (0.7243 - 0.009)/0.702 = 1.019$~; $\Phi =
0.8459$~: le 85e centile.
\textbf{19.} $422 \times (1 - 0.8459) = 65$.
\textbf{20.} $M = T_0\,e^{\mu + s^2/2}\,\Phi(z + s) = 1200 \times 0.00442 \times 0.9570 =
\qty{5.08}{Mt}$.
\textbf{21.} $1200 \times 0.00442 \times 0.5203 = \qty{2.76}{Mt}$.
\textbf{22.} $\Phi(z + s) = 0.8364$ du cuivre, pour $\Phi(z) = 0.6663$ du \omterm{def:b2:ore-deposits:grade}{tonnage}.
\textbf{23.} De 0.3 à \qty{0.2}{\percent}, le \omterm{def:b2:ore-deposits:grade}{tonnage} croît de $1054 - 800 = \qty{254}{Mt}$
(\qty{32}{\percent}) et le cuivre de $5.08 - 4.44 = \qty{0.64}{Mt}$ (\qty{14}{\percent}) : les blocs ajoutés
ont une moyenne de $0.64/254 = \qty{0.25}{\percent}$, les plus pauvres du \omterm{def:b2:ore-deposits:deposit}{gisement}.
\textbf{24.} $1200 \times 0.00442 \times 0.8364 = 800 \times 0.00555$~: l'indice contient
\textbf{\qty[text-series-to-math=true]{4.44}{Mt}} de cuivre au-dessus de la coupure de \qty{0.3}{\percent}.
\end{solution}
